\chapter*{\appendixname}%

\section{Primer apéndice}%
\label{apen:apen1}%
Los apéndices son contenido esencial pero complementario: incluyen material elaborado por los autores de la tesis, que es relevante para entender la investigación, pero demasiado extenso o técnico para el cuerpo principal.\par%

\section{Segundo apéndice}%
\label{apen:apen2}%
Algunos ejemplos de apéndices:%
\begin{itemize}
    \item códigos fuente completos,%
    \item demostraciones matemáticas extensas,%
    \item cuestionarios utilizados en la metodología.%
\end{itemize}
 
Los ápendices guardan relación directa con el trabajo: son parte integral de la investigación y se citan en el texto principal.\par%